% REFERÊNCIAS — Volume 2 (gerado por expandir_volume.py)

\chapter{Referências Bibliográficas}

\begin{itemize}[leftmargin=1.6em]
  \item BRASIL. Ministério da Educação. Base Nacional Comum Curricular. Brasília: MEC, 2018.
  \item BRASIL. Ministério da Educação. Política Nacional de Alfabetização (Decreto nº 9.765, de 11 de abril de 2019). Brasília: MEC, 2019.
  \item KUMON, Toru. Kumon: o segredo da mãe japonesa. Rio de Janeiro: Campus, 1996.
  \item DEHAENE, Stanislas. Os neurônios da leitura: como a ciência explica a nossa capacidade de ler. Porto Alegre: Penso, 2012.
  \item SOARES, Magda. Alfaletrar: toda criança pode aprender a ler e a escrever. São Paulo: Contexto, 2020.
  \item MORAIS, Artur Gomes de. Sistema de escrita alfabética. São Paulo: Melhoramentos, 2012.
  \item SALLES, Jerusa F. et al. Manual de avaliação neuropsicológica infantil. São Paulo: Vetor, 2016.
  \item CAPOVILLA, Alessandra G. S.; CAPOVILLA, Fernando C. Alfabetização: método fônico. São Paulo: Memnon, 2009.
\end{itemize}

\chapter{Tabela de Correspondência Fonema-Grafema Completa}

A tabela lista os fonemas do português brasileiro com suas grafias principais e as grafias trabalhadas neste volume.

\begin{center}
\renewcommand{\arraystretch}{1.5}
\begin{tabular}{lll}
\toprule
**Fonema** & **Grafias** & **Exemplos** \\
\midrule

/ch/ & ch & chave, chuva \\

/lh/ & lh & ilha, coelho \\

/nh/ & nh & ninho, banho \\

/ku/ & qu & queijo, quilo \\

/gu/ & gu & guerra, água \\

/br/ & br & braço, livro \\

/kr/ & cr & cravo, preto \\

/dr/ & dr & dragão, madre \\

/fr/ & fr & fruta, frio \\

/gr/ & gr & grade, grande \\

/pr/ & pr & prato, primo \\

/tr/ & tr & trem, corte \\

/vr/ & vr & livro, palavra \\

/r/ (forte) & rr & carro, terra \\

/s/ (forte) & ss & passar, osso \\

/s/, /z/ & s/z & casa, zero \\

/s/ (ç) & ç & coração, açúcar \\

\bottomrule\end{tabular}\end{center}

\chapter{Glossário de Termos Técnicos}

\begin{description}
  \item[Consciência fonológica] capacidade de refletir sobre os sons da fala, das sílabas aos fonemas.
  \item[Dígrafo] duas letras que representam um único fonema (ch, lh, nh, rr, ss, qu, gu).
  \item[Encontro consonantal] duas consoantes que ocorrem juntas na sílaba, cada uma com seu som (br, cr, bl, cl).
  \item[Fluência] leitura com precisão, ritmo adequado e expressividade, que libera recursos para a compreensão.
  \item[Fonema] menor unidade sonora da língua; o som de CH em chave é um fonema.
  \item[Grafema] representação escrita de um fonema; o dígrafo ch é um grafema de duas letras.
  \item[Nasalidade] passagem do ar pelo nariz na produção de sons (m, n, nh, ÃO, ÃE, ÕE).
  \item[Automatismo] leitura e escrita sem esforço consciente, obtido por treino distribuído e repetitivo.
  \item[Rastreio escolar] observação pedagógica sistemática e sequenciada que orienta o apoio; nunca diagnostica.
  \item[Rotas 1A-1D] percursos internos de progressão usados pelo professor para diferenciar o ensino por evidência.
\end{description}

\chapter{Bibliografia Comentada}

Cada obra abaixo fundamenta uma peça do método. O professor pode ler a obra inteira ou o capítulo indicado para aprofundar.

\begin{description}

  \item[KUMON, Toru. Kumon: o segredo da mãe japonesa. Rio de Janeiro: Campus, 1996.] Apresenta a filosofia do treino diário, da gradação mínima e da repetição até o automatismo.

  \item[DEHAENE, Stanislas. Os neurônios da leitura. Porto Alegre: Penso, 2012.] Explica a via fonológica da leitura e por que a decodificação automatizada precede a compreensão.

  \item[SOARES, Magda. Alfaletrar: toda criança pode aprender a ler e a escrever. São Paulo: Contexto, 2020.] Sintetiza a didática da alfabetização e a passagem do alfabético ao ortográfico.

  \item[MORAIS, Artur Gomes de. Sistema de escrita alfabética. São Paulo: Melhoramentos, 2012.] Apoia a exploração das regularidades ortográficas do português, como as trabalhadas no 2º ano.

  \item[CAPOVILLA, Alessandra; CAPOVILLA, Fernando. Alfabetização: método fônico. São Paulo: Memnon, 2009.] Descreve a sequência fônica e os exercícios de consciência fonológica que inspiram a Parte VII.

  \item[SALLES, Jerusa F. et al. Manual de avaliação neuropsicológica infantil. São Paulo: Vetor, 2016.] Fundamenta a observação de leitura, escrita e funções executivas em contexto escolar.

  \item[BRASIL. Base Nacional Comum Curricular. Brasília: MEC, 2018.] Documento oficial que define as habilidades de Língua Portuguesa por ano.

  \item[BRASIL. Política Nacional de Alfabetização (Decreto nº 9.765/2019). Brasília: MEC, 2019.] Orienta a ênfase em evidências científicas e no método fônico na alfabetização.

  \item[CAGLIARI, Luiz Carlos. Alfabetização e linguística. São Paulo: Scipione, 1989.] Discute a consciência fonológica na aquisição da escrita e o papel da variação linguística.

  \item[FERREIRO, Emilia; TEBEROSKY, Ana. Psicogênese da língua escrita. Porto Alegre: Artmed, 1999.] Descreve as hipóteses de escrita que fundamentam a observação diagnóstica do professor.

\end{description}

\chapter{Integridade e Ética no Uso do Livro}

Este material é de uso pedagógico e reproduzível para fins escolares não comerciais. O rastreio integrado é observação pedagógica e nunca diagnóstico; a triagem não substitui a avaliação do profissional habilitado. Nenhum dado individual deve ser compartilhado sem autorização da família, conforme a LGPD (Lei nº 13.709/2018). O professor que aplicar as fichas do Volume Profissional deve ter a formação exigida pelo instrumento.
Ao reproduzir as planchetas, preserve o cabeçalho com o nome, a data e o descritor; ao arquivar as folhas, use o Registro de Progresso individual, nunca comparações públicas entre alunos.

\chapter{Conectividade com a Rede de Ensino}

Este volume se integra ao currículo da rede: os descritores CAEd e as matrizes do SPAECE estão mapeados nas unidades e avaliações; o planejamento bimestral conversa com o calendário escolar; o Registro de Progresso alimenta os conselhos de classe. A família recebe, na reunião de pais, o resumo do método e o combinado de 10 minutos diários de leitura em casa.
